\question[5]
Ein Geheimtext wurde mit dem Vigenère-Verfahren verschlüsselt.
Die Schlüssellänge wurde bereits als $3$ bestimmt.

Der Geheimtext lautet:

\begin{center}
\texttt{OIGORXBKXX}
\end{center}

Zur weiteren Analyse werden die Buchstaben entsprechend ihrer Position aufgeteilt:

\begin{center}
\begin{tabular}{|c|c|c|}
\hline
1. Teiltext & 2. Teiltext & 3. Teiltext \\
\hline
O O B X & I R K & G X X \\
\hline
\end{tabular}
\end{center}

Im ersten Teiltext tritt der Buchstabe \texttt{O} am häufigsten auf.

\begin{enumerate}[label=\alph*)]
\item Erkläre, warum der Geheimtext in Teiltexte zerlegt wird.
\item Nimm an, dass \texttt{O} für den häufigen deutschen Buchstaben \texttt{E} steht. Bestimme die zugehörige Caesar-Verschiebung für den ersten Teiltext.
\end{enumerate}

\begin{solutionbox}{5cm}
\begin{enumerate}[label=\alph*)]
\item Alle Buchstaben eines Teiltexts wurden mit derselben Verschiebung verschlüsselt. Jeder Teiltext kann daher wie ein Caesar-Geheimtext untersucht werden.
\item E $\rightarrow$ O entspricht einer Verschiebung um 10 Stellen.
Die vermutete Caesar-Verschiebung lautet daher $10$.
\end{enumerate}
\end{solutionbox}